% answers_short.tex —— 2006 年·普通高等学校招生考试·湖南卷文科 · 简短答案速查
% 由 tools/make_short_answers.py 自 answers.tex 的 \gksolhead/\gkanswer 提取，
% 并与 2006/2006_数学真题_湖南文_答案要点.md 逐题比对后生成，勿手改（改源 MD 后重跑）。

\begin{gknotice}
\textbf{使用说明}：本节只给最终答案便于快速核对，完整推导见解析版。答案为编写者独立推导，\textbf{非官方标准答案}。
\end{gknotice}

\gkpart{选择题}
\begin{center}\begin{tabular}{ccccccccccc}
\toprule
题号 & 1 & 2 & 3 & 4 & 5 & 6 & 7 & 8 & 9 & 10 \\
\midrule
答案 & D & C & D & A & A & B & C & B & D & C \\
\bottomrule
\end{tabular}\end{center}

\gkpart{填空题}
\begin{itemize}[leftmargin=2.6em]
  \item[\textbf{11.}] $2^n-1$。
  \item[\textbf{12.}] 85。
  \item[\textbf{13.}] 5。
  \item[\textbf{14.}] 6 条。
  \item[\textbf{15.}] $a=-3$。
\end{itemize}

\gkpart{解答题}
\begin{itemize}[leftmargin=2.6em]
  \item[\textbf{16.}] $\theta=\dfrac\pi3$ 或 $\dfrac{2\pi}3$。
  \item[\textbf{17.}] (1) $0.31$；(2) $0.90$；(3) U/PENDING
  \item[\textbf{18.}] (1) $PQ\perp$ 平面 $ABCD$；(2) $\arccos\dfrac13$；(3) $2\sqrt2$
  \item[\textbf{19.}] (1) 当 $a>0$ 时，递增区间为 $(-\infty,0)$、$(2/a,+\infty)$，递减区间为 $(0,2/a)$；当 $a<0$ 时，递减区间为 $(-\infty,2/a)$、$(0,+\infty)$，递增区间为 $(2/a,0)$；$a=0$ 不在定义域内；(2) $a\in[-1,0)\cup[3,4]$
  \item[\textbf{20.}] (1) $a_4=10,\ a_5=15,\ a_n=\dfrac{n(n+1)}2$；(2) $b_n=2+\dfrac4{n(n+2)}$，$b_1+b_2+\cdots+b_n=2n+3-\dfrac2{n+1}-\dfrac2{n+2}$
  \item[\textbf{21.}] (1) $m=0,\ p=\dfrac98$；抛物线焦点为 $(\dfrac9{16},0)$，不在 $AB:x=1$ 上；(2) $m=\pm\dfrac{\sqrt6}{3}$；对应 $AB$ 方程为 $y=-3m(x-1)$
\end{itemize}
